\section{Generated torsors}
\label{generated-torsors}
Fix a group~$G$. Given a set~$X$ and a map~$\gamma:X\to G$, we say that $X$ \emph{generates}~$G$ (with respect to~$\gamma$) when $\freegroup\gamma:\freegroup X\to G$ is surjective.
% (in particular, given a presentation $P=\pres XR$, $X$ always generates the presented group~$\pgroup P$ through the canonical morphism).
From now on, we suppose that we are in such a situation.
%
We now provide a variant for the construction of a delooping of~$G$ by $G$-torsors described in the
previous section, taking advantage of the additional data of a generating
set in order to obtain smaller and simpler constructions.
%
Note that here, contrarily to \cref{delooping-hit}, we only need a set of
generators, not a full presentation.

\paragraph{Actions of sets} 
Given a type~$A$, we write $\End A$ for its type of \emph{endomorphisms}, \ie
maps $A\to A$. An \emph{action} of the set~$X$ on a set~$A$ is a morphism
$X\to\End A$, \ie a family of endomorphisms of~$A$ indexed by~$X$.
%
We write $\XSet$ for the type
\[
  \XSet
  \qdefd
  \Sigma(A:\Set).(X\to\End A)
  \]
% $\Sigma(A:\Set).(X\to\End A)$
of actions of~$X$. An element~$\alpha$ of this type consists in a set $\dom\alpha$ with a function
$\alpha:X\to\End(\dom\alpha)$ and is called an \emph{$X$-set}. A
\emph{morphism} between $X$-sets $\alpha$ and $\beta$ is a function
$f:\dom\alpha\to\dom\beta$ satisfying \cref{G-set-morphism} for every $x:X$.
%
The identities between $X$-sets can be characterized in a similar way as for
$G$-sets, see \cref{Gset-equalities}, and \cref{GSet-SIP} also extends in the
expected way.

Precomposition by $\gamma$ induces a function $U:\GSet\to\XSet$ which can be
thought of as a forgetful functor from $G$-sets to $X$-sets. Note that $U$
depends on~$\gamma$ but we leave it implicit for concision.

\paragraph{Applications of the generated delooping}
We have seen in the previous section that the connected
component of the principal $G$-torsor~$P_G$ in $G$-sets is a delooping
of~$G$. Our aim in this section is to show here that this construction can be
simplified by taking the connected component of the restriction of $P_G$ to $X$-sets.

Before proving this theorem, which is formally stated as
\cref{generated-delooping} below, we shall first illustrate its use on a
concrete example. Consider $\Z_n$, the cyclic group with $n$ elements. We write
$s:\Z_n\to\Z_n$ for the successor (modulo $n$) function, which is an
isomorphism. By \cref{PG-delooping}, we know that the type
\[
  \Sigma(A:\GSet[\Z_n]).\ptrunc{\P{\Z_n}=A}
\]
of $\Z_n$-torsors is a model of~$\B\Z_n$.
%
This type is the connected component of the principal $\Z_n$-torsor $\P{\Z_n }$
in the universe~$\GSet[\Z_n]$ of sets with an action of $\Z_n$, \ie
sets~$A$ equipped with a morphism $ \alpha:\Z_n\to\Aut A $. Such a set~$A$ is
thus comes with one automorphism $\alpha(k)$ for every element~$k:\Z_n$,
therefore $k$ automorphisms in this case. However, most of them are superfluous:
$1$ generates all the elements of~$\Z_n$ by addition, so $\alpha(1)$
generates all the $\alpha(k)$ by composition because $\alpha(k)=\alpha(1)^k$.
%
The useful data of a $\Z_n$-set thus boils down to a set~$A$ together with one
automorphism $\alpha:\Aut A$ such that $\alpha^n=\id_A$.

Indeed, writing
$
  \SetEnd
  \defd
  \Sigma(A:\Set).\End A
$
for the type of all \emph{endomorphisms} (on any set), our theorem will imply
that the type
\begin{equation}
  \label{Zn-delooping}
  \Sigma((A,f):\SetEnd).\ptrunc{(\Z_n,s)=(A,f)}
\end{equation}
(the connected component of the successor modulo $n$ in the universe of set
endomorphisms) is still a delooping of $\Z_n$. Note that we didn't assume
that~$f$ is an isomorphism nor that it should verify $f^n=\id$. This is because
both properties follow from the fact that $f$ is in the connected component of
the successor (which satisfies those properties). Similarly, we do not need to
explicitly assume that the domain of the endomorphism is a set.

Our theorem thus allows to define, in a relatively simple way, types
corresponding to deloopings of groups. As recalled in the introduction, this is
particularly useful when one is not disposed to use higher inductive types
(\eg because their definition, implementation and semantics are not entirely
mature). This is in fact the reason why this approach was used in UnitMath to
define the circle~\cite{bezem2021construction}, and we provide a generic way to
similarly define other types.
%
We expect that it can be used to reason about groups and compute invariants such as their cohomology~\cite{cavallo2015synthetic,buchholtz2020cellular,brunerie2022synthetic}.
%
On a side note, one might be worried by the fact that we are ``biased'' (by
using a particular set of generators), which allows us to be more concise but
might make more difficult generic proofs compared to $G$-torsors: we expect that
this is not the case because in order to define the group~$G$ itself, one
usually needs to resort to a presentation, and thus is also biased in some sense...

\paragraph{The generated delooping}
In the following, we write $\P X$ for $U\P G$.

\begin{proposition}[\agda{XSetProperties}{theorem}]
  \label{PG-loop-restriction}
  We have a group equivalence
  $
  \Loop\P G\equivto\Loop \P X
  $.
\end{proposition}
\begin{proof}
  From \cref{Gset-equalities}, an element of $\Loop P_G$ consists of an equality
  $p:G=G$ in $\U$ such that
  \[
    \P G(x)\circ\transport p=\transport p\circ\P G(x)
  \]
  for every $x:G$. By function extensionality and the definition of the action $\P G$, this is equivalent to
  requiring, for every $g,z:G$ that
  \begin{equation}
    \label{eq:PG-condition}
    g(\transport p(z))=\transport p(gz)
  \end{equation}
  Note that the above equality is between elements of~$G$, which is a set, and is
  thus a proposition.
  %
  Similarly, an element of $\Loop \P X$ consists of an equality $p:G=G$
  in~$\U$ satisfying
  \begin{equation}
    \label{eq:PX-condition}
    \gamma(x)(\transport p(z))=\transport p(\gamma(x)z)
  \end{equation}
  for every $x:X$ and $z:G$.

  Clearly, any equality $p:G=G$ in $\Loop\P G$ also belongs to $\Loop \P X$
  since the condition \cref{eq:PX-condition} is a particular case of
  \cref{eq:PG-condition}. We thus have a function
  $\phi:\Loop\P G\to\Loop \P X$.
  %
  Conversely, consider an element $p:G=G$ of $\Loop \P X$, \ie satisfying
  \cref{eq:PX-condition} for every $x:X$ and $z:G$. Our aim is to show that it
  belongs to $\Loop\P G$. Given $g,z:G$, we thus want to show that
  \cref{eq:PG-condition} holds. Since~$\freegroup{\gamma}$ is surjective,
  because $X$ generates $G$, we know that there merely
  exists an element $u$ of~$\freegroup{X}$ such that
  $\freegroup{\gamma}(u)=x$. Since \cref{eq:PG-condition} is a proposition, by
  the elimination principle of propositional truncation, we can actually suppose
  given such a $u$, and we have
  \begin{align*}
    x(\transport p(y))
    &=\freegroup\gamma(u)(\transport p(y))&&\text{since $\freegroup{\gamma}(u)=x$}\\
    &=\transport p(\freegroup\gamma(u)y)&&\text{by repeated application of \cref{eq:PX-condition}}\\
    &=\transport p(xy)&&\text{since $\freegroup{\gamma}(u)=x$.}
  \end{align*}
  The second equality essentially corresponds to the commutation of the
  following diagram, where $u\defd x_1x_2\ldots x_n$ with $x_i:X$:
  \[
    \begin{tikzcd}
      G
      \ar[d,"\transport p"']
      % \ar[rrrrr,bend left,"\freegroup\gamma(u)"]
      \ar[r,"\gamma(x_1)"]
      &
      G
      \ar[d,"\transport p"']
      \ar[r,"\gamma(x_2)"]
      &
      G
      \ar[d,"\transport p"']
      \ar[r]
      &
      \ldots
      \ar[r]
      &
      G
      \ar[d,"\transport p"]
      \ar[r,"\gamma(x_n)"]
      &
      G
      \ar[d,"\transport p"]
      \\
      G
      \ar[r,"\gamma(x_1)"']
      &
      G
      \ar[r,"\gamma(x_2)"']
      &
      G
      \ar[r]
      &
      \ldots
      \ar[r]
      &
      G
      \ar[r,"\gamma(x_n)"']
      &
      G
    \end{tikzcd}
  \]
  % can be detailed as follows. Writing $u=x_1\ldots x_n$ with $x_i:X$, we have:
  % \EO{Faire une preuve diagrammatique !!}  \SM{en fait je remplacerais ça par 3
    % cas: $u$ est un générateur $a$, $u$ est un produit et $u$ est une unité...}
  % \begin{align*}
    % \freegroup\gamma(u)(\transport p(y))
    % &=\freegroup\gamma(x_1\ldots x_n)(\transport p(y))
    % &&\text{by definition of~$u$}
    % \\
    % &=\freegroup\gamma(x_1\ldots x_{n-1})\gamma(x_n)(\transport p(y))
    % &&\text{$\freegroup\gamma$ is a group morphism extending $\gamma$}
    % \\
    % &=\freegroup\gamma(x_1\ldots x_{n-1})(\transport p(\gamma(x_n)y))
    % &&\text{by \cref{eq:PX-condition}}
    % \\
    % &=\freegroup\gamma(x_1\ldots x_{n-2})\gamma(x_{n-1})(\transport p(\gamma(x_n)y))
    % &&\text{$\freegroup\gamma$ is a group morphism extending $\gamma$}
    % \\
    % &=\freegroup\gamma(x_1\ldots x_{n-2})(\transport p(\gamma(x_{n-1})\gamma(x_n)y))
    % &&\text{by \cref{eq:PX-condition}}
    % \\
    % &=\freegroup\gamma(x_1\ldots x_{n-2})(\transport p(\freegroup\gamma(x_{n-1}x_n)y))
    % &&\text{$\freegroup\gamma$ is a group morphism extending $\gamma$}
    % \\
    % &\vdots
    % \\
    % &=\freegroup\gamma(1)(\transport p(\freegroup\gamma(x_1\ldots x_n)y))
    % \\
    % &=1(\transport p(\freegroup\gamma(x_1\ldots x_n)y))
    % &&\text{$\freegroup\gamma$ is a group morphism}
    % \\
    % &=\transport p(\freegroup\gamma(x_1\ldots x_n)y)
    % &&\text{$1$ is a unit}
    % \\
    % &=\transport p(\freegroup\gamma(u)y)
    % &&\text{by definition of~$u$}
  % \end{align*}
  This thus induces a function $\psi:\Loop \P X\to\Loop\P G$. The functions
  $\phi$ and $\psi$ clearly preserve the group structure (given by
  concatenation of paths) and are mutually inverse of each other, hence we have the equivalence we wanted.
\end{proof}

\begin{theorem}
  \label{generated-delooping}
  The type $\Comp \P X$ is a delooping of $G$.
\end{theorem}
\begin{proof}
  Since taking the connected component preserves loops spaces (\cref{B-loop}),
  we have that $\Loop\Comp\P X$ is equal to $\Loop\P X$, which in turn is equal
  to $\Loop\P G$ by \cref{PG-loop-restriction}, and thus to $G$ by
  \cref{PG-delooping}.
\end{proof}

The delooping of~$G$ constructed in previous theorem is the component of~$\P X$
in $X$-sets:
\[
  \Sigma(A:\U).\Sigma(S:\isSet A).\Sigma(f:X\to\End A).\ptrunc{\P X=(A,S,f)}
\]
Since for any type~$A$, the type $\isSet A$ is a
proposition~\cite[Theorem~7.1.10]{hottbook}, and the underlying type of~$\P X$
is a set, the underlying type of any $X$-set in the connected component
of~$\P X$ will also be a set. As a consequence, the above type can slightly be
simplified, by dropping the requirement that~$A$ should be a set:

\begin{proposition}
  The type $\Sigma(A:\U).\Sigma(f:X \to\End A).\ptrunc{\P X=(A,f)}$ is a delooping
  of~$G$.
\end{proposition}

\noindent
For instance, the delooping \cref{Zn-delooping} of~$\Z_n$ can slightly be
simplified as
\[
  \Sigma((A,f):\Endomorphisms).\ptrunc{(\Z_n,s)=(A,f)}
\]
where $\Endomorphisms\defd\Sigma(A:\U).\End A$ is the type of all endomorphisms
of the universe.

\begin{example}
  \label{Dn}
  \Cref{generated-delooping} applies to every group for which a generating set
  is known (and, of course, the smaller the better). For instance, given a
  natural number~$n$, the \emph{dihedral group}~$D_n$ is the group of symmetries
  of a regular polygon with~$n$ sides. It has $2n$ elements and is generated by
  two elements $s$ (axial symmetry) and $r$ (rotation by an angle of
  $2\pi/n$). Hence the connected component of the symmetry and the rotation in
  the type of pairs of set endomorphisms, \ie
  \[
    \Sigma(A:\Set).\Sigma(f, g : \End A\times\End A).\ptrunc{(D_n,s,r)=(A,f,g)}
  \]
  is a delooping of the dihedral group~$D_n$.
\end{example}

\paragraph{Alternative proof}
We would like to provide another proof \cref{PG-loop-restriction}, which was suggested by an anonymous reviewer.
It is based on the idea that in order to show $\Loop\P X\equivto\Loop\P G$, it is enough to show that $U:\GSet\to\XSet$ is an \emph{embedding}, \ie that for every $\alpha,\beta:\GSet$ the induced function $U^=:(\alpha=\beta)\to(U\alpha=U\beta)$ between path spaces is an equivalence~\cite[Definition~4.6.1]{hottbook}.
It relies on the following result whose proof can be found in \cref{omitted-proofs}.

\begin{lemma}
  \label{fiberwise-embedding}
  Given a type~$A:\U$, type families $P,Q:A\to\U$ and $f:(a:A)\to P\,a\to Q\,a$, the map $\Sigma A.f:\Sigma A.P\to\Sigma A.Q$ is an embedding if and only
  if $f\,a:P\,a\to Q\,a$ is an embedding for every $a:A$.
\end{lemma}

\begin{proposition}
  The map $U:\GSet\to\XSet$ is an embedding.
\end{proposition}
\begin{proof}
Given a morphism of groups $f:H\to K$, we write $\GSet[f]:\GSet[K]\to\GSet[H]$
for the function induced by precomposition. In particular, by definition, we
have $U\defd\GSet[\gamma]$. The function $\gamma:X\to G$ can be decomposed as
$\freegroup\gamma\circ\incl$, and therefore $\GSet[\gamma]$ can be decomposed as
$\GSet[\incl]\circ\GSet[\freegroup\gamma]$:
\[
  \begin{tikzcd}
    &\GSet[\freegroup X]\ar[dr,"{\GSet[\incl]}"]&\\
    \GSet\ar[ur,"{\GSet[\freegroup\gamma]}"]\ar[rr,"U"']&&\XSet
  \end{tikzcd}
\]
Both maps are embeddings so that $U$ is an embedding by composition. Namely,
\begin{itemize}
\item \emph{$\GSet[\freegroup\gamma]$ is an embedding}. Consider the map
  \[
    F:(A:\Set)\to(G\toGroup\Aut A)\to(\freegroup X\toGroup\Aut A)
  \]
  obtained by precomposition by $\freegroup\gamma$. Since $\freegroup\gamma$ is
  surjective we have that $F\,A$ is an
  embedding for every set $A$~\cite[Lemma~10.1.4]{hottbook}. By \cref{fiberwise-embedding}, we
  deduce that $\Sigma\Set.F$, which is $\GSet[\freegroup\gamma]$, is an
  embedding.
\item \emph{$\GSet[\incl]$ is an embedding}. By universal property of $\freegroup X$, given a set~$A$, the map
  \[
    (\freegroup X\toGroup\Aut A)\to(X\to\Aut A)
  \]
  obtained by precomposition by $\incl$ is an equivalence, and thus an
  embedding. Moreover, the property of being an isomorphism is a proposition, hence the forgetful map
  \[
    (X\to\Aut A)\to(X\to\End A)
  \]
  is an embedding because its fibers are propositions. Since embeddings are
  stable under composition, we deduce that the induced map
  \[
    F:(A:\Set)\to(\freegroup X\toGroup\Aut A)\to(X\to\End A)
  \]
  is such that $F\,A$ is an embedding for every set $A$. By
  \cref{fiberwise-embedding}, the map $\Sigma\Set.F$, which is $\GSet[\incl]$, is
  thus an embedding.\qedhere
\end{itemize}  
\end{proof}

% Local Variables:
% TeX-master: "article"
% End:
